\section{Testing Strategy}
To ensure the reliability of the application, a multi-phased testing approach was adopted:
\begin{itemize}
    \item \textbf{Functional Testing}: Each feature (e.g., creating a formation, registering a participant) was tested against its functional requirements to ensure correctness.
    \item \textbf{Validation Checks}: Input fields were tested for edge cases, including empty values, incorrect formats (e.g., invalid emails), and numeric range violations.
    \item \textbf{API Testing}: Backend endpoints were validated using tools like Postman to verify status codes, JSON response structures, and security constraints.
\end{itemize}

\section{Results and Discussion}
The development of TrainingMS successfully transitioned the organization's training management from a fragmented manual process to a centralized digital platform. 

\subsection{Improvements Over Manual System}
The primary achievements include:
\begin{itemize}
    \item \textbf{Data Accuracy}: Centralized storage eliminated the data duplication issues inherent in yearly Excel files.
    \item \textbf{Efficiency}: Search and retrieval times for participant records and session history were reduced from minutes to seconds.
    \item \textbf{Visibility}: The statistics module provided insights into training trends that were previously impossible to generate manually.
\end{itemize}

\subsection{Strengths and Limitations}
The application's strengths lie in its modular architecture and user-centric design. However, some limitations remain, such as the current lack of a mobile-specific interface and the absence of an automated email notification system for participants. These areas provide opportunities for future enhancements.
